\section{引子}

本节介绍两个例子，以说明假设检验在实际生产生活中的作用。

~

例1. 某面粉厂用自动包装机包装面粉，正常情况下一袋面粉的重量（克）呈正态分布$N\left( 500,9 \right) $。
现随机抽取30袋质检，检测其重量是否达标，发现本次抽样到的平均重量为498.8g，样本均值比总体期望少了1.2g。
此时，我们要判断缺少的这1.2g是抽样造成的，还是包装机的问题，若是后者，则须考虑停机检修。

~

例2. 抛硬币100次，统计正反面朝上的次数如下。
按道理抛100次正反面几乎差不多，62和38相差有点悬殊，我们需要判断硬币是不是有问题。
\begin{table}[h]
\centering
\begin{tabular}{ccc}
    \toprule
    向上的面 & 正面 & 反面\\
    \midrule
    次数 & 62 & 38\\
    \bottomrule
\end{tabular}
\end{table}

~

上面的问题都可以归结到概率问题。
例1可以描述成容量30的样本的均值出现1.2偏差的概率，例2可以描述成p=0.5的伯努利试验出现62-38的概率。
若概率小，说明本次试验居然发生了小概率事件，而抽样和统计没问题，剩下的问题最大可能就出在总体上了。




